\documentclass[10pt]{article}
\usepackage[T1]{fontenc}
\usepackage[utf8]{inputenc}
\usepackage[french]{babel}
\usepackage{amsmath,amssymb,mathrsfs,mathpazo}
\usepackage{geometry}
\geometry{a4paper,margin=2cm}
\usepackage{graphicx,array,tabularx,tabularray,colortbl,multirow,multicol,booktabs,enumitem,pifont,ulem,fancybox,diagbox,cancel}
\usepackage{pstricks,pstricks-add,pst-plot,pst-tree,pst-node,pst-text,pst-func,pst-math,pst-bspline,pst-3dplot,pst-eucl,pst-solides3d}
\let\pscalculate\undefined
\usepackage{pst-bezier}
\usepackage{tikz,pgfplots,tkz-tab}
\usetikzlibrary{arrows,arrows.meta,patterns,calc,positioning,trees,decorations.markings,shapes.geometric,shapes.misc,shapes.arrows,intersections,decorations.pathmorphing,decorations.pathreplacing,backgrounds,shadows,babel}
\pgfplotsset{compat=1.18}
\usepgfplotslibrary{fillbetween}
\usepackage[np]{numprint}
\usepackage{listings,makecell,pas-tableur,twoopt}
\let\barre\undefined
\usepackage[pstricks]{tablvar}
\definecolor{darkgreen}{RGB}{0,100,0}
\providecommand{\red}{\color{red}}
\providecommand{\blue}{\color{blue}}
\providecommand{\green}{\color{green}}
\providecommand{\black}{\color{black}}
\providecommand{\white}{\color{white}}
\providecommand{\gray}{\color{gray}}
\providecommand{\cyan}{\color{cyan}}
\providecommand{\magenta}{\color{magenta}}
\providecommand{\yellow}{\color{yellow}}
\providecommand{\violet}{\color{violet}}
\providecommand{\orange}{\color{orange}}
\providecommand{\euro}{\texteuro}
\providecommand{\EUR}{\texteuro}
\providecommand{\ieme}{\textsuperscript{e}}
\providecommand{\degres}{\ensuremath{^\circ}}
\providecommand{\degree}{\ensuremath{^\circ}}
\providecommand{\vv}[1]{\overrightarrow{#1}}
\providecommand{\Coord}[1]{#1}
\providecommand{\decosix}{}
\providecommand{\starredbullet}{\ensuremath{\bullet}}
\pagestyle{empty}
\setlength{\parindent}{0pt}
\usepackage{adjustbox}\DeclareUnicodeCharacter{2236}{:}\DeclareUnicodeCharacter{2190}{\ensuremath{\leftarrow}}
\providecommand{\R}{\mathbb{R}}
\providecommand{\N}{\mathbb{N}}
\providecommand{\Z}{\mathbb{Z}}
\providecommand{\Q}{\mathbb{Q}}
\providecommand{\D}{\mathbb{D}}
\providecommand{\C}{\mathbb{C}}
\providecommand{\e}{{\mathrm{e}}}
\providecommand{\vect}[1]{\overrightarrow{#1}}
\providecommand{\vectt}[1]{\overrightarrow{\mathrm{#1}}}
\providecommand{\vvt}[1]{\overrightarrow{\mathrm{#1}}}
\providecommand{\np}[2][]{{#2~#1}}
\providecommand{\vv}[1]{{\overrightarrow{#1}}}
\renewcommand{\d}{{\mathrm{d}}}
\providecommand{\degres}{{^\circ}}
\providecommand{\degre}{{^\circ}}
\providecommand{\blue}{\color{blue}}
\providecommand{\red}{\color{red}}
\providecommand{\enskip}{\quad}
\providecommand{\llbracket}{\lbrack\!\lbrack}
\providecommand{\rrbracket}{\rbrack\!\rbrack}
\providecommand{\euro}{{\text{€}}}
\providecommand{\no}{{\text{n°}}}
\providecommand{\interval}[2]{{[#1;#2]}}
\providecommand{\intervFO}[2]{{[#1;#2[}}
\providecommand{\dotfill}{\ldots}

\providecommand{\ds}{\displaystyle}
\providecommand{\textding}[1]{#1}
\providecommand{\Oij}{\(\mathrm{O};\vec{\imath},\vec{\jmath}\)}
\providecommand{\Oijk}{\(\mathrm{O};\vec{\imath},\vec{\jmath},\vec{k}\)}
\providecommand{\Ouv}{\(\mathrm{O};\vec{u},\vec{v}\)}
\providecommand{\intervFF}[2]{\ensuremath{[#1;#2]}}
\renewcommand{\euro}{\ifmmode\text{\texteuro}\else\texteuro\fi}
\DeclareUnicodeCharacter{00B0}{\ifmmode^{\circ}\else\textdegree\fi}
\DeclareUnicodeCharacter{20AC}{\ifmmode\text{\texteuro}\else\texteuro\fi}
% Énoncés : Éducation nationale. Transcription et corrigés : APMEP, Denis Vergès et contributeurs. Découpage et classement : CoopMaths, Sébastien Lozano, Éric Elter et contributeurs. Adaptation de la présentation : Maths Explorer.
% CC BY-SA : https://coopmaths.fr/www/about/
% https://coopmaths.fr/alea/static/bac/2022/tex/bac_2022_05_sujet2_metropole_2.tex
\begin{document}
\noindent\textit{APMEP / CoopMaths. CC BY-SA. Présentation adaptée par Maths Explorer.}\par\medskip
\medskip

\emph{Cet exercice est un questionnaire à choix multiples. Pour chacune des questions suivantes, une seule des quatre réponses proposées est exacte. Une réponse fausse, une réponse multiple ou l'absence de réponse à une question ne rapporte ni n'enlève de point.\\
Pour répondre, indiquer sur la copie le numéro de la question et la lettre de la réponse choisie.\\Aucune justification n'est demandée.}

\medskip

Pour les questions 1 à 3 ci-dessous, on considère une fonction $f$ définie et deux fois dérivable sur $\R$. La courbe de sa fonction dérivée $f'$ est donnée ci-dessous.

On admet que $f'$ admet un maximum en $- \dfrac{3}{2}$ et que sa courbe coupe l'axe des abscisses au point de coordonnées $\left(- \dfrac12~;~0\right)$.

\begin{minipage}{0.5\linewidth}
\textbf{Question 1 }:

\textbf{a.~~} La fonction $f$ admet un maximum en $- \dfrac{3}{2}$ ;

\textbf{b.~~}  La fonction $f$ admet un maximum en $- \dfrac{1}{2}$ ;

\textbf{c.~~} La fonction $f$ admet un minimum en $- \dfrac{1}{2}$; 

\textbf{d.~~}  Au point d'abscisse $-1$, la courbe de la
fonction $f$ admet une tangente horizontale.
\end{minipage}\hfill
\begin{minipage}{0.48\linewidth}
\begin{center}
On rappelle que la courbe ci-dessous représente la fonction dérivée $f'$ de $f$.
\end{center}

\psset{unit=1.2cm}
\begin{pspicture*}(-5.2,-2.6)(1,1)
\psgrid[gridlabels=0pt,subgriddiv=4,gridwidth=0.25pt,subgridwidth=0.15pt]
\psaxes[linewidth=1.25pt,labelFontSize=\scriptstyle]{->}(0,0)(-5.2,-2.6)(1,1)
\psplot[plotpoints=2000,linewidth=1.25pt,linecolor=red]{-5}{3}{2 x mul 1 add 2.71828 x exp mul neg}
\end{pspicture*}
\end{minipage}

\medskip

\textbf{Question 2 }:

\begin{center}
\begin{tabularx}{\linewidth}{X X}
\textbf{a.~~}La fonction $f$ est convexe sur $\left]- \infty~;~- \dfrac32\right[$ ;&
\textbf{b.~~}La fonction $f$ est convexe sur $\left]- \infty~;~- \dfrac12\right[$ ;\\
\textbf{c.~~}La courbe $\mathcal{C}_f$ représentant la fonction $f$ n'admet pas de point d'inflexion ; &\textbf{d.~~}La fonction $f$ est concave sur $\left] - \infty~;~- \dfrac12\right[$.
\end{tabularx}
\end{center}

\medskip

\textbf{Question 3}:

La dérivée seconde $f''$ de la fonction $f$ vérifie :

\begin{center}
\begin{tabularx}{\linewidth}{X X}
\textbf{a.~~} $f''(x) \geqslant  0$ pour $x \in \left]-\infty~;~- \dfrac12\right[$ ; &\textbf{b.~~} $f''(x) \geqslant  0$ pour $x \in [- 2~;~- 1]$ ;\\
\textbf{c.~~} $f''\left(- \dfrac32 \right) = 0$ ;&\textbf{d.~~} $f''(- 3) = 0$.
\end{tabularx}
\end{center}


\textbf{Question 4 :}

\medskip

On considère trois suites $\left(u_n\right)$, $\left(v_n\right)$ et $\left(w_n\right)$. On sait que, pour tout entier naturel $n$, on a : $u_n \leqslant v_n\leqslant  w_n$ et de plus: $\displaystyle\lim_{n \to + \infty} u_n= 1$ et $\displaystyle\lim_{n \to + \infty} w_n= 3$.

On peut alors affirmer que :

\begin{center}
\begin{tabularx}{\linewidth}{X X}
\textbf{a.~~} la suite $\left(v_n\right)$ converge ;&\textbf{b.~~} Si la suite 
$\left(u_n\right)$ est croissante alors la suite $\left(v_n\right)$ est minorée par $u_0$ ;\\
\textbf{c.~~} $1 \leqslant  v_0 \leqslant 3$ ;&\textbf{d.~~} la suite $\left(v_n\right)$ diverge.
\end{tabularx}
\end{center}


\medskip

\textbf{Question 5}:

\medskip

On considère une suite $\left(u_n\right)$ telle que, pour tout entier naturel $n$ non nul: $u_n \leqslant  u_{n+1}  \leqslant \dfrac1n$.

On peut alors affirmer que : 

\begin{center}
\begin{tabularx}{\linewidth}{X X}
\textbf{a.~~}la suite $\left(u_n\right)$ diverge ;&\textbf{b.~~}la suite $\left(u_n\right)$ converge ;\\
\textbf{c.~~}$\displaystyle\lim_{n \to + \infty} u_n =  0$ ;&\textbf{d.~~} $\displaystyle\lim_{n \to + \infty} u_n =  1$.
\end{tabularx}
\end{center}

\medskip

\textbf{Question 6}:

\medskip

On considère $\left(u_n\right)$ une suite réelle telle que pour tout entier naturel $n$, on a : $n < u_n < n + 1$. 

On peut affirmer que:

\begin{center}
\begin{tabularx}{\linewidth}{X X}
\textbf{a.~~}Il existe un entier naturel $N$ tel que $u_N$ est un entier ;&\textbf{b.~~} la suite $\left(u_n\right)$ est croissante ;\\
\textbf{c.~~} la suite $\left(u_n\right)$ est convergente ;&\textbf{d.~~}La suite $\left(u_n\right)$ n'a pas de limite.
\end{tabularx}
\end{center}



\bigskip

\end{document}
